\section{Method}
\label{sec:method}

\subsection{One interception point}
Every LLM call and tool call of a workflow goes through a single runtime object. Workflow code asks
for \texttt{llm(node, messages)}, \texttt{tool(name, args)} or \texttt{agent(node, prompt, tools)}
(a standard tool-calling loop); a recipe tells the runtime what to do with each request: answer it
from a recording, fake it with an LLM, merge it with its neighbours into one call, or run it for
real. The workflow's own code never changes, so every miniature is the same program as the full
workflow with different call semantics, and every run leaves a trace (calls, sources, tokens,
cost, latency) from which the judge and the analysis work.

\subsection{Reduction operators}
Table~\ref{tab:operators} lists the operators, grouped by how much validity they put at risk.

\begin{table*}[t]
\centering\small
\begin{tabular}{@{}llp{6.3cm}p{5.4cm}@{}}
\toprule
& \textbf{Operator} & \textbf{What it does} & \textbf{Risk to validity} \\
\midrule
\multirow{3}{*}{A} & \op{A1} prefix replay & Calls before the focus have identical inputs to a recorded run of the base variant; answer them from the recording. & None: the inputs are identical. \\
& \op{A2} tool cache & Replay recorded real tool outputs for identical calls. & Stale outputs for stateful tools. \\
& \op{A3} paired runs & Variants share identical calls (and simulated outputs). & None; reduces variance. \\
\midrule
\multirow{4}{*}{B} & \op{B1} slice & Drop nodes whose output never reaches what the judge scores. & Hidden shared state. \\
& \op{B2} cap loops & Bound agent steps or conversation turns. & Variants that win through persistence. \\
& \op{B3} cheaper model off-focus & Run the nodes not under test on a cheaper model. & Changes the focus's inputs. \\
& \op{B4} fewer tasks & A fixed random subset of the tasks. & Sampling noise. \\
\midrule
\multirow{6}{*}{C} & \op{C1} simulated tools & An LLM writes each tool's output: from the description (\emph{pure}), with recorded examples from other tasks (\emph{grounded}), or only for calls never recorded (\emph{cache-first}). & Too clean (no errors, no mess); no state. \\
& \op{C2} call fusion & One call writes a whole agent loop (or a whole pipeline) as a transcript with imagined tool results. & The model imagines results that suit its plan. \\
& \op{C3} prompt shrinking & Compress the fixed part of a prompt. & Behaviour shifts silently. \\
& \op{C4} sub-agent stub & A non-focus sub-agent becomes one fused call. & Misses what real search finds. \\
& \op{C5} early stop & Stop after the node whose output the judge needs. & Drops a later correction. \\
& \op{C6} short observations & Cut every tool output to a fixed length. & Facts beyond the cut. \\
\bottomrule
\end{tabular}
\caption{Reduction operators. Group A cannot change any decision; group B rarely does; group C must
be checked.}
\label{tab:operators}
\end{table*}

Two design choices keep the simulated parts honest. First, a tool simulator never sees the user's
question, only the tool's description and arguments (and, when grounded, real outputs of the same
tool): a real tool does not know the question either, and a simulator that does is an easy-mode
oracle. Second, simulated outputs are memoised per call, so all variants see the same fake world
and differences between them come from the variants.

\subsection{Choosing operators}
For each part of a workflow, six properties decide which operators are safe: its \emph{distance}
from the focus (far parts can be shrunk hard), its \emph{cost share}, whether it is
\emph{deterministic} (replay it), whether its output \emph{reaches the judge} (else slice it), how
\emph{fakeable} it is (a calculator is, a stateful database is not), and whether the focus
\emph{reacts} to its output (then a fake must reproduce failures too). In short: keep real what you
change, replay what you did not change, simulate what you cannot replay, and check before you
trust.

\subsection{Scoring miniatures of stateful benchmarks}
Some benchmarks score the world a run leaves behind; $\tau$-bench \citep{yao2024tau} compares the
final database with the one the expected actions produce. A miniature whose tools are simulated
changes no database, so we score every run, full or miniature, by replaying the agent's
\emph{write} calls from its trace on a fresh copy of the real database. For real runs this equals
$\tau$-bench's own check; for miniatures it scores the agent's decisions rather than the
simulator's side effects.

\subsection{Judging}
Clear-cut runs are scored by rule (normalised exact match, or the benchmark's own check). All
others go to a blind LLM judge, a fresh Claude Code instance with no tools, given batches of runs
under random codes: the task, the reference, the final answer and a short log. It never sees the
recipe, the variant, or whether a run came from a miniature. For every failure it returns one
reason from a fixed per-workflow taxonomy and one sentence of explanation. When a benchmark has
its own pass/fail check, the judge supplies only the reason.
